% ==================== 第 3 章 系统设计与实现 ====================
\section{系统设计与实现}

\subsection{系统总体架构}

\begin{figure}[H]\centering
\includegraphics[width=0.96\textwidth]{\detokenize{system_architrcture.png}}
\caption{系统总体架构：LangGraph 引擎与质量流水线 V3.4、外部服务与数据源、资产与存储层}\label{fig:arch}
\end{figure}

\noindent 一次查询的七环节闭环流程展示见图~\ref{fig:reallflow}（第 1 章）。

\subsubsection{模块总表}

\begin{longtable}{@{}p{1.7cm}p{2.1cm}p{6.2cm}p{2.4cm}p{2.2cm}@{}}
\toprule
实际模块 & 输入 & 本作品的核心处理 & 输出 & 与下一模块的关系 \\
\midrule
\endfirsthead
\toprule
实际模块 & 输入 & 本作品的核心处理 & 输出 & 与下一模块的关系 \\
\midrule
\endhead
\textbf{任务理解} & 自然语言查询（如“昴星团的年龄、距离和金属丰度”）与用户上传的论文 & 先识别查询类型（天文查询、寒暄、退出意图、无效输入），寒暄或无效输入给出友好回应并引导；对天文查询以多轮对话逐项澄清目标天体与所需性质，最后汇总确认；随后借助天文数据库确认查询主体标准名称，结合领域知识库筛选目标性质，生成统一的性质清单 & 目标天体标识、性质清单（字段名、单位、含义）与澄清轮次记录 & 性质清单是后续检索、提取与校验环节共同遵循的字段约定 \\
\textbf{数据检索} & 性质清单与天体别名 & 并行匹配多个专业星表；按性质生成文献检索式并排序；自动下载开放获取论文；识别论文附带的数据表 & 候选来源列表与本地论文库 & 各来源以统一编号关联，零产出的来源在聚合时过滤 \\
\textbf{数据提取} & 论文各页图像 & 用多模态大模型从正文和表格中提取数值，数值与单位严格分离并受性质清单约束；为每条数据定位原文位置；从论文图表中筛选相关图证 & 结构化的数值记录（含原文摘录与位置）、图表证据、原文定位页面 & 记录标明来源类型（正文、表格或星表），供后续环节区分处理 \\
\textbf{质量检查} & 提取出的全部记录 & 对每条记录逐条评估：字段正确性（类型、单位写法、格式、合理范围）、来源可回溯性（页码、原文定位框、原文摘录）、提取置信度；对同一物理量的多来源取值做方差分析，发现异常与冲突组 & 问题清单、逐源可靠性得分、异常与冲突组 & 问题清单驱动下一步清洗 \\
\textbf{数据清洗} & 问题清单与待清洗记录 & 先按问题类型执行清洗：统一数值格式与单位换算、去除重复或近似重复记录、补齐缺失项、将列名与写法对齐到标准字段；随后对多源取值做冲突分类与标注，异常值送人工复核。清洗与冲突处理在此环节内部循环迭代，最多三轮，全程修改留痕 & 清洗后的记录、冲突标注与修改记录 & 清洗结果经校验后进入交付环节 \\
\textbf{数据交付} & 清洗后的记录 & 将数据导出为长表、宽表与结构化文件，附质量报告、来源信息与修改记录 & 最终结构化数据、质量报告、来源信息与修改记录 & 输出供前端展示与下载，作为洞察环节的输入 \\
\textbf{数据洞察} & 最终结构化数据 & 生成字段关联分析、跨字段关系、使用建议与综合叙述报告，全部明确标注为模型生成 & 洞察报告 & 与结构化数据一并呈现给研究者，供核验与使用 \\
\bottomrule
\end{longtable}

\subsubsection{数据闭环的关键设计}

\begin{itemize}[leftmargin=1.2em,itemsep=2pt]
\item \textbf{为什么不是一次性检索或内容摘要：}普通单次检索或大模型摘要往往依赖单一来源，一旦出现提取偏差或单位混淆便直接导致结论失效。本作品面向真实科研场景，强调“多源海量获取 + 交叉印证”，逐条提取数值、物理量纲与精准证据位置，从根本上解决单点数据不可靠与无法验证的问题。
\item \textbf{质量问题如何形成修正闭环（闭环机制与设计逻辑）：}
\begin{enumerate}[leftmargin=1.6em,itemsep=1pt]
\item \textbf{前置多源冗余，消除对“外部重搜”的被动依赖}：系统在检索阶段通过专业星表、论文正文及补充材料构建高覆盖度的数据候选池（单次查询通常汇聚多篇文献与几十条记录）。这种广覆盖的来源冗余，使得系统无需在单条数据出错时进行盲目、高成本的外部重搜。
\item \textbf{聚焦内层质量筛选与规则自愈}：系统的修正闭环集中在数据可用性提升上：
\begin{itemize}[leftmargin=1.2em,itemsep=1pt]
\item \textbf{格式与量纲自愈}：在对齐环节对不对称误差、括号不确定度及单位形态混乱进行自动消歧与两段式换算，若规则执行异常则由状态机就近重新规划；
\item \textbf{真假冲突消解}：在整合环节发现数值分歧时，自动将数据打回对齐层统一物理口径，消除因量纲不一致带来的伪冲突。
\end{itemize}
\item \textbf{严格质量门禁与冲突保真}：对最终未达标的数据坚决执行隔离与剔除，杜绝低质数据流入下游；对有效的历史测量分歧（如不同巡天时代的系统偏差）进行保留并显式归因。最终依托剩余的高置信度多源记录相互印证，交付兼具高可用性与学术参考价值的结构化成果。
\end{enumerate}
\end{itemize}

\subsection{Qwen 与上下文工程}

\subsubsection{Qwen 在本作品中的实际作用}

\begin{longtable}{@{}p{2.4cm}p{2.6cm}p{4.6cm}p{4.9cm}@{}}
\toprule
使用场景 & 使用的模型 & 提供给模型的上下文 & 输出与约束 \\
\midrule
\endfirsthead
\toprule
使用场景 & 使用的模型 & 提供给模型的上下文 & 输出与约束 \\
\midrule
\endhead
意图澄清与寒暄处理 & qwen3.7-flash & 用户查询文本与已澄清的历史轮次 & 识别查询类型（天文查询、寒暄、退出、无效）并解析目标天体与所需性质；寒暄与无效输入给出引导回应 \\
目标性质筛选 & qwen3.8-max & 天文数据库解析出的天体标识与类型、领域知识库中的候选性质 & 从候选性质中筛选与查询匹配的性质，生成统一的性质清单（字段名、单位、含义），字段名唯一 \\
文献检索式生成 & qwen3.7-flash & 目标天体别名与所需性质清单 & 按性质分组生成文献检索式，交给检索模块执行 \\
论文数值提取 & qwen3.7-plus 多模态模型 & 论文各页图像、性质清单作为提取白名单、页码映射说明、输出示例 & 从正文与表格提取数值记录（数值与单位分离、与原文逐字一致、附原文摘录），字段名必须取自性质清单，置信度必填 \\
原文定位兜底 & qwen3.7-flash 多模态模型 & 高清裁剪的局部页面图、数值候选写法与上下文片段 & 对本地定位失败的数值给出定位框，或对表格给出整表定位框；坐标经校验后映射回整页 \\
图证提取与相关性判定 & qwen3.7-flash 多模态模型 & 论文页面图、目标天体与所需性质 & 检测页面中的图表区域并判定与查询的相关性，输出定位框、图注与判定理由；表格、公式明确排除 \\
质量评估 & qwen3.7-flash & 记录数据与领域评价规则 & 逐条评估字段正确性、来源可回溯性与提取置信度，输出问题清单与逐源可靠性得分 \\
清洗方案生成 & qwen3.7-flash & 该来源的问题清单、示例与失败反馈 & 生成数据清洗方案（格式清理、数值规整等），交由本地规则执行并复核；禁止生成单位换算代码，越界修改被护栏拦截 \\
冲突分类 & qwen3.7-flash & 同一物理量的多来源取值统计 & 对取值分群给出冲突类型与置信标注，供研究者参考 \\
洞察报告生成 & qwen3.7-flash & 最终结构化数据与领域知识库 & 生成字段关联分析、跨字段关系、使用建议与综合叙述，全部明确标注为模型生成 \\
结果摘要与前端呈现 & qwen3.7-flash & 任务运行结果与各环节统计 & 生成任务摘要与前端展示文本，供研究者快速了解查询结果 \\
模型与调用方式 & 全部经阿里云百炼兼容接口调用 & 结构化输出统一约束（响应格式 + 多级解析回退），输出经结构校验后才进入下游 & 模型不直接访问外部资源，检索由本地模块执行，定位由本地定位器完成，清洗由本地规则执行 \\
\bottomrule
\end{longtable}

\noindent\emph{（Qwen 调用凭证：阿里云百炼控制台调用日志与计费凭证截图随提交材料附上）}

\subsubsection{上下文组织与约束控制}

\begin{itemize}[leftmargin=1.2em,itemsep=2pt]
\item \textbf{团队为减少无来源生成（幻觉）、字段错配或错误合并实际采取的措施：}① 提取前注入性质清单作白名单，模型只能从中选择字段；② 要求数值与单位严格分离、与原文逐字一致，提供原文上下文规范思考链，禁止自行换算或改写；③ 提取置信度低于阈值的记录直接丢弃；④ 数值的原文位置以本地逆向定位为准——从论文页面文本层按数值与上下文反查真实位置，而非直接采信模型给出的定位框，防止定位幻觉；⑤ 洞察与建议类输出一律明确标注为模型生成，与原始数据区分。
\item \textbf{上下文在数据查找、解析和修正前后如何更新：}检索前注入天体标识与性质清单；提取时按论文逐篇注入页码映射；清洗时注入评估发现的问题清单与上一轮失败反馈；每轮循环的结果（修改记录）回填上下文，供后续决策参考；任务完成后上下文不参与下游，只保留结构化结果与修改记录。
\item \textbf{该设计对数据结果带来的实际变化：}提取字段全部落在性质清单内，未再出现清单外字段；数值与单位分离约束使“数值内嵌单位”类脏格式显著减少；原文定位框使每条论文记录都可点击回溯到原始页面。
\end{itemize}

\noindent\emph{（上下文结构示意图待补充：研究目标、数据需求、字段字典、候选来源、原始片段、历史结果与质量反馈如何进入模型）}

\subsection{检索与提取}

（对应模板 P8 与 P9 设问；M45 实况卡贯穿本节）

\begin{figure}[H]\centering
\includegraphics[width=0.9\textwidth]{\detokenize{search_retrieve.png}}
\caption{检索三路并行与多模态提取链：星表/文献/补充材料 → PropertySpec 字段白名单 → VLM 提取}\label{fig:retrieval}
\end{figure}

\subsubsection{检索与来源筛选}

\noindent\textbf{查找过程：}

\begin{longtable}{@{}p{2.3cm}p{6.8cm}p{2.8cm}p{2.6cm}@{}}
\toprule
查找环节 & 本作品实际做法 & 中间结果 & 对后续处理的作用 \\
\midrule
确定查找目标 & 从澄清确认的目标天体与性质清单出发，明确需要检索的物理量、目标别名与来源类型 & 查找目标清单 & 决定星表匹配键与文献检索式的方向 \\
统一天体标识 & 在天文数据库中解析目标，取得标准名称、天体类型与完整别名清单 & 统一的天体标识 & 作为星表匹配与文献检索的共同依据，避免别名漏检 \\
并行多源检索 & 星表按别名与坐标匹配查询；文献库按分组检索式检索，检索式由模型按性质生成 & 星表命中结果与论文候选列表 & 提供两类主要来源的记录 \\
获取论文全文 & 按论文编号多级尝试下载开放获取全文，无法获取的论文标记为不可用 & 本地论文库 & 决定哪些论文可进入数值提取 \\
补充来源识别 & 识别论文附带的数据表，作为正文提取的补充 & 补充数据表清单 & 提供正文未报告的原始数值 \\
来源筛选与去重 & 按检索相关度排序，评估来源可靠性，去除重复来源 & 筛选后的来源清单 & 保证进入解析环节的来源可用且相关 \\
缺失处理与收尾 & 无法获取全文的来源尝试其他渠道后仍失败则零产出并过滤；全部来源处理完毕后停止 & 最终来源清单 & 决定最终输出的来源覆盖 \\
\bottomrule
\end{longtable}

\noindent\textbf{来源可用性判断：}

\begin{itemize}[leftmargin=1.2em,itemsep=2pt]
\item \textbf{团队如何判断某一来源可以用于本研究需求：}① 来源相关性，星表按天体别名匹配星表筛选，文献按照检索相关度排序；② 来源可靠性，评估阶段为每个来源计算可靠性得分，低分来源不进入清洗环节；③ 字段相关性，提取前以统一的性质清单约束，无对应性质的内容放弃而非硬塞进其他字段。
\item \textbf{来源中无法核实或不适合纳入的数据如何处理：}无法核实的来源（封闭期刊无开放全文）直接零产出并在输出中过滤；不适合的内容（公式、统计量、参考推导值）在提取阶段以提示词的方式明确排除；无法确认是否为同一对象的数据送人工复核。
\end{itemize}

\noindent\textbf{M45 检索实况卡。}以真实查询昴星团（M45）的年龄、距离和金属丰度为例：确定查找目标后，在天文数据库中将 M45 解析为标准名称与别名清单（含 Pleiades、Cl Melotte 22 等），随后并行执行两类检索：

\begin{itemize}[leftmargin=1.2em,itemsep=1pt]
\item \textbf{星表检索}：在专业星表库中按星团类型与别名匹配，命中 MWSC 星团目录，取得星团年龄、距离、距离模数、金属丰度等 4 条目录记录。
\item \textbf{文献检索}：在文献库中按模型生成的三组性质分组检索式逐一检索（检索式针对昴星团别名与对应物理量关键词组合，如 \code{(title:"Pleiades" OR abs:"Pleiades") AND (abs:age OR abs:ages OR abs:"logarithmic age" OR abs:"log t")}，距离组、金属丰度组同构）。检索结果按相关度排序，取得候选论文列表后，逐篇尝试经开放获取渠道下载全文——最终成功下载 29 篇（1998--2025 年，10 种期刊），全部进入数值提取。筛选时按检索相关度与来源可靠性排序，保留报告了目标性质且取得全文的论文；排除无法取得开放全文的文献。
\end{itemize}

\noindent 候选来源精选（下表为代表性精选，完整 30 个来源清单随附录交付）：

\begin{longtable}{@{}p{3.4cm}p{1.8cm}p{4.2cm}p{1.6cm}p{3.4cm}@{}}
\toprule
候选来源 & 来源类型 & 找到的相关内容 & 是否采用 & 采用或排除的实际理由 \\
\midrule
MWSC 星团目录（J/A+A/558/A53） & 数据库 & 年龄、距离、距离模数、金属丰度 4 条目录记录 & 是 & 权威星团目录，字段覆盖查询要求 \\
2023A\&A...677A.163A（Gaia 天体测量视角） & 论文 & 视差、距离 135.15 $\pm$ 0.43 pc、成员星分析 & 是 & 基于 Gaia 的现代测距结果，检索相关度高 \\
2005A\&A...429.887P（昴星团距离） & 论文 & 距离 133.8 $\pm$ 3 pc（主序拟合） & 是 & 直接报告目标性质的测距研究 \\
2017A\&A...598A.48G（移动星团法修订距离） & 论文 & 距离 134.4 +2.9 $-$2.8 pc（运动学方法） & 是 & 不同测距方法交叉印证 \\
1999ApJ...523.328N（Hipparcos 视差误差） & 论文 & 距离 130.7 $\pm$ 11.1 pc（早期测距） & 是 & 历史测距值，供冲突保留对比 \\
2016A\&A...595A..59M（恒星双胞胎法距离） & 论文 & 距离 134.8 $\pm$ 1.7 pc & 是 & 现代测距结果，多源互证 \\
\bottomrule
\end{longtable}

\subsubsection{提取与解析}

\noindent\textbf{实际处理的来源类型：}

\begin{longtable}{@{}p{2.2cm}p{5.2cm}p{2.8cm}p{2.2cm}p{2.2cm}@{}}
\toprule
实际来源类型 & 使用的解析方法 & 提取的内容 & 质量检查方式 & 无法解析时如何处理 \\
\midrule
数据库（星表） & 按天体别名匹配专业星表，配合表列源数据过滤所需性质及数值，直接取得目录记录 & 星表的年龄、距离、星等、金属丰度、视差等目录值 & 来源可靠性评分、字段完整性检查、查询结果去重 & 查询失败自动切换备用镜像，仍失败则该星表零产出 \\
论文正文 & 将论文各页渲染为图像，用多模态大模型逐篇提取正文中的数值，要求数值与单位分离、附原文摘录；随后在页面文本层按数值与上下文逆向定位真实位置 & 正文中报告的目标性质数值、单位、原文摘录、页码与定位框 & 白名单字段校验、置信度阈值、数值与原文逐字一致核对、定位结果校验 & 提取失败自动重试，重试无果该论文降级；定位失败降级为段落级定位，不丢弃数据 \\
论文表格 & 多模态大模型识别表格结构并提取单元格数值，单位从表头或上下文判定；随后按表格标题框出整表区域作为溯源范围 & 表格中的数值记录、单位、表标题、整表定位框 & 数值与单位分离校验、表内取值合理范围检查、整表定位包含性校验 & 表格结构无法识别时该表零产出；定位失败降级为标题区域定位 \\
论文图证 & 多模态模型检测页面中的图表区域并判断与查询的相关性，相关图按定位框裁剪保存 & 论文图表证据（裁剪图）、图注、相关性说明 & 无，依靠提示词约束质量 & 页面无相关图时该页零产出；图表定位失败时该图不保存 \\
论文补充材料 & 按论文编号识别其附带的数据表，对表结构与列语义做判定，确认与查询相关才纳入 & 补充数据表中正文未报告的原始数值 & 与目标实体对应性检查 & 无法判定表内容与查询相关时整表丢弃，不纳入提取 \\
\bottomrule
\end{longtable}

\noindent\textbf{M45 提取实况卡：一项复杂来源的解析示例。}以真实查询昴星团（M45）为例，五种来源各自提取到的数据如下：

\begin{itemize}[leftmargin=1.2em,itemsep=1pt]
\item \textbf{数据库（MWSC 星团目录）}：经结构化查询按星团类型与别名匹配，取得 4 条目录记录——年龄 8.15（对数写法）、距离模数 5.576 mag、距离 130 pc、金属丰度 $-$0.036（对数写法），均保留原始单位写法并附数据表行列位置。
\item \textbf{论文正文}：如 1999ApJ 第 13 页正文“Hipparcos 视差误差相关性”上下文处提取距离 130.7 $\pm$ 11.1 pc，数值与单位分离、附原文摘录与页码，并经页面文本层逆向定位到该数值所在行；2016A\&A 第 7 页正文提取距离 134.8 $\pm$ 1.7 pc，同法处理，共提取 64 条正文原始记录。
\item \textbf{论文表格}：如 2005A\&A 第 1 页表格提取金属丰度 $-$0.034 $\pm$ 0.024（单位 dex），表头上下文判定单位；2006A\&A 第 2 页“星团参数汇总表”表格提取年龄 75-100 Myr 与金属丰度 $-$0.14 to +0.13。共提取 35 条表格原始记录。
\item \textbf{论文图证}：从 29 篇论文中筛选出 78 张相关图表（如 1999ApJ 第 20 页 Fig.~5，展示 Hipparcos 视差的平滑处理结果，直接呈现平均视差与距离模数），每张附裁剪图、图注与相关性判定理由。
\item \textbf{论文补充材料}：本案例对 29 篇论文的附带数据表逐一识别与列语义判定，均未能确认与查询直接相关，按保守策略整表丢弃，未纳入任何记录（该机制在历史查询中曾成功阻断整表污染源）。
\end{itemize}

\subsection{质量管线（自动质量机制）}

\begin{figure}[H]\centering
\includegraphics[width=0.9\textwidth]{\detokenize{quality.png}}
\caption{质量管线全景：评估/规范化/流程控制/人工审核/冲突检测/数据导出/洞察}\label{fig:quality}
\end{figure}

\noindent 分工说明：本部分讲\textbf{系统运行时自动}的质量评估、清洗、冲突与整合交付机制；\textbf{人工}发现质量问题并修复的记录见第 4 章。本节不再重复第 4 章内容。

\subsubsection{质量管线架构与自动评估机制}

\noindent\textbf{整体流向}：提取结果进入质量管线后依次经过 Assessment 子图的四段评估（\textbf{Profiling → QualityAssessment → Scoring → Decision}）→ 按来源\textbf{派发}到处理路径（Normalization / Conflict / HumanReview / Export）→ 规范化与冲突经 LoopController 受控循环（规范化⇄冲突最多 3 轮，全局最多 8 次迭代，超限强制导出，保证流程有界收敛）→ 六段\textbf{导出} → 洞察 → 结束。质量管线的输出驱动后续一切：评估的问题清单驱动清洗、方差分析驱动冲突、冲突组与低质量来源驱动人工复核（见本章图 4 与 3.4.2--3.4.4 实况）。

\bigskip
\noindent\textbf{评估环节（纯规则与统计优先，模型为辅）}：Profiling 阶段为\textbf{零 LLM 的纯统计预检}，先摸清各来源记录面貌；QualityAssessment 对每个来源逐条评估 \textbf{5 个维度}——提取质量 extraction\_quality、完整性 completeness、一致性 consistency、格式合规 format、来源可靠性 source\_reliability，并对同一物理量的跨源取值做全局方差分析（analyze\_multi\_source\_variance）；Scoring 将 5 维 + 冲突风险 conflict\_risk 共 \textbf{6 键按领域权重加权}得到逐源得分与综合质量分（overall，前端 ×100 展示），并按记录量校准出置信区间（0.85 这类综合分即该加权结果）；Decision 为纯规则判据：提取分过低 / 跨实体标识错误 / 越界记录过半 → 直接人工审核；检出问题 → 规范化；检出异常组 → 冲突分析；仅有方差或无问题 → 导出。另以“质量 × 修复代价”矩阵做单向升级（严重度 Export $<$ Normalization $<$ Conflict $<$ HumanReview），保证严重问题只升不降。

\bigskip
\noindent\textbf{四路处理路径与分工：}

\begin{longtable}{@{}p{2.5cm}p{3.4cm}p{5.8cm}p{2.6cm}@{}}
\toprule
处理路径 & 触发条件 & 执行内容 & 出口 \\
\midrule
规范化 Normalization & 评估问题清单（缺单位/格式问题/字段别名等，条件→工具 9 条映射）或冲突消解方案要求 & SourceRouter → Planning → Normalization → Validation → Report 五链；坚持\textbf{模型生成方案 + 本地规则/沙箱执行}分离（L1 确定性工具、L2 规则注入、L3 模板沙箱运行 Python，先 dry-run 5 条再全量）；校验失败回退 Planning（≤2 次） & 通过 → 冲突复查或导出；否则回流 \\
冲突分析 Conflict & 方差分析检出明显分群 / 异常 & 按（实体、字段）聚合：方差聚合 → 差异分类 → 异常验证 → 置信度评估 → 消解报告；用 Cohen's d（$>$2.0 判统计离群）、年差（$>$5 年判时域差异）、Bootstrap KS 检验等统计量区分\textbf{方法差异、条件差异、时域演化、测量不确定度、重复观测}等 6 类原因并给验证置信度；unit\_error 类自动转单位归一，cross\_id\_error 类送人工 & 消解报告驱动规范化或人工复核；全部取值保留（preserve\_all） \\
人工审核 HumanReview & 提取质量极低、跨实体标识错误、越界过半、冲突无法自动裁决 & 批量协议一次挂起全部待审项，逐项 \textbf{1--5 裁决}（1 采用来源 A / 2 采用来源 B / 3 自定义值 / 4 双值保留标注 / 5 跳过），可附理由，理由写入修改记录与导出溯源；纯质量审核项仅给 4/5 选项；无效输入重问 ≤3 次后按跳过落定 & 提交 → 规范化执行修改；无法判断 → 回到评估重估；取消 → 数据原样交付 \\
导出 Export & 其余路径完成后 & 六段链：组织 → Schema 格式化 → 元数据 → 溯源构建 → 输出校验 → 结构化文件生成；校验失败仍组装但标记隔离（quarantine）不入正式结果 & 五件套：结构化数据 / 长表 / 宽表 / 质量总结 / 元数据 + 溯源 + 洞察报告写盘 \\
\bottomrule
\end{longtable}

\noindent\textbf{全程留痕}：规范化与冲突的每一次修改都写入 modifications 统计与 data\_trace 记录（字段、修改前值、修改后值、执行 agent、工具、原因、置信度），修改前后数据均可恢复——这是“保留修正前后真实差异”与前端“处理轨迹”展示的数据基础。

\bigskip
\noindent 对应模板 P5 的“团队评价口径”层面（评哪些方面为何 / 评价主体 / 五类状态定义 / 修正前后口径一致）在第 4 章 4.1 声明。M45 案例的自动评估产出（综合分 0.85、逐源得分、26 项异常、31 项冲突归因）见 3.4.3/3.4.4 实况。

\subsubsection{数据统一：把不同写法归一}

对应模板 P10。本节目的是\textbf{表示层统一}：把“同一个意思的不同写法”翻译成同一套标准表示（字段名、对象标识、数值类型与格式、单位、缺失标记）。凡能按规则换算或映射的差异在此处理；凡需要裁决的竞争取值与记录取舍见 3.4.3（P11）。去重属记录层决策，已移至 3.4.3。

\noindent\textbf{统一过程：}

\begin{longtable}{@{}p{2.2cm}p{3.4cm}p{5.4cm}p{2.4cm}p{1.9cm}@{}}
\toprule
需要统一的内容 & 不同来源实际怎么写 & 本作品实际怎么处理 & 统一后的结果 & 仍无法确认的内容 \\
\midrule
字段名称与含义 & 星表列名与论文表述各异（如“距太阳的距离”“d”与“距离”） & 在任务理解阶段用标准性质库生成统一的性质清单，提取时字段名必须取自该清单；残余异构写法在清洗环节按字段映射归一到标准字段 & 所有来源的记录归入同一标准字段体系 & 无法可靠映射的写法保留原始名称并标注 \\
对象、材料或记录的名称与标识 & 同一目标有多个别名（如 M45、Pleiades、Cl Melotte 22） & 以天文数据库解析的标准名称与天体类型为身份锚点，别名清单用于各来源匹配 & 不同别名归并为同一目标 & 无法确认是否为同一目标时送人工复核 \\
数据类型 & 数值以字符串、科学计数法、带前缀或范围等形式出现 & 清洗环节将字符串数值统一为数值类型，清理前缀与多余格式 & 数值格式统一，可参与后续计算 & 区间或带不确定度的写法保留原样 \\
单位 & 同一量在不同来源以不同单位呈现（如 Myr、yr、log(yr)） & 清洗环节按配置的换算规则统一单位（含对数写法还原），换算全程留痕 & 单位统一到标准写法 & 换算无法可靠判定时保留原始写法 \\
缺失内容和特殊状态 & 部分来源缺少某字段，或出现空值 & 保留缺失标记，不填充 & 缺失与特殊状态被显式标识 & 无法判断的数据送人工复核 \\
\bottomrule
\end{longtable}

\noindent\textbf{M45 清洗实况卡。}以真实查询昴星团（M45）为例，对照上表各统一维度，质量管线中的实际处理如下：

\begin{itemize}[leftmargin=1.2em,itemsep=1pt]
\item \textbf{字段名称与含义}：MWSC 星团目录的原始列名 logt、d、MOD、[Fe/H] 经统一归入标准字段 age、distance、dist\_modulus、fe\_h，原始列名保留在来源信息中供核对。
\item \textbf{对象名称与标识}：查询文本中的“昴星团（M45）”经天文数据库解析，与别名 Pleiades、Cl Melotte 22 归并为同一目标（标准名称 Cl Melotte 22），各来源记录同属该目标。
\item \textbf{数据类型}：清洗环节将字符串数值统一为数值类型，真实修改记录如 “134.4 +2.9 $-$2.8” → “134.4”（提取中心值）、“7.48 $\pm$ 0.03” → “7.48”（清理不确定度）、“~ 120-130” → “120-130”（清理前缀）、“5.60” → “5.6”（数值规整），共 69 条修改全程留痕。
\item \textbf{单位}：昴星团（M45）查询中，星团年龄出现 Myr、yr、log(yr) 三种写法，清洗环节按性质库标准单位将 7 条 Myr 记录换算为 yr；星团目录 MWSC 的对数写法 log(yr)（年龄）与 log(Sun)（金属丰度）因对数形态无法可靠互算，保留原始写法并标注；距离（pc）、视差（mas）、距离模数（mag）各来源写法一致，无需换算。
\item \textbf{缺失内容和特殊状态}：另一真实查询（北河三 Pollux）中，四个星表均未提供查询要求的质量、半径字段（只提供温度与光谱型），系统保留缺失标记不填充，缺失字段在质量报告中明确列出。
\end{itemize}

\subsubsection{整合、冲突与来源保留：记录层决策}

对应模板 P11。P10（3.4.2）把“写法”归一后，本节处理\textbf{记录层的整合决策}：面对同一对象的多条记录与竞争取值，如何取舍、裁决与保真——凡不能靠换算解决的差异在此裁决。

\noindent\textbf{整合情况表：}

\begin{longtable}{@{}p{2.3cm}p{3.6cm}p{4.8cm}p{3.3cm}@{}}
\toprule
整合情况 & 本作品如何判断 & 实际处理方式 & 输出如何保留来源或差异 \\
\midrule
同一对象的多条记录 & 以天文数据库解析的标准名称与天体类型为身份锚点，各来源记录归属同一目标 & 按来源分别保留为独立记录，不合并 & 每条记录保留各自来源标识与原文位置 \\
重复或近似重复记录 & 清洗环节检测同来源内的重复或近似重复条目 & 去除重复记录，去除行为留痕 & 保留的记录保留完整来源信息 \\
不同来源数值冲突 & 评估环节对同一物理量的多来源取值做方差分析，识别明显分群 & 全部取值保留，标注冲突类型与值域区间；无法裁决的送人工复核 & 冲突值逐条保留，标注中写明涉及来源与取值区间 \\
部分字段缺失 & 评估环节检查各来源预期应含字段是否缺失 & 保留缺失标记，不填充 & 缺失字段在质量报告中明确列出 \\
无法确认是否同一对象 & 实体归属存疑或跨来源实体错配异常 & 送人工复核，由研究者裁决 & 以待审核状态呈现，不自动合并 \\
\bottomrule
\end{longtable}

\noindent 聚合的物理执行位置：检索子图末端 result\_aggregator 节点把三路命中（SIMBAD 天体解析、星表查询、论文检索）汇合为统一编号的输出包并记录零产出来源（即 3.1 模块表中“聚合”所在）；此后记录进入质量管线评估。上表第 1 行“同一对象多条记录按来源独立保留”的取舍决策由质量管线执行，并在导出时以独立记录携带各自来源信息。

\bigskip
\noindent\textbf{M45 冲突实况卡（真实运行记录，任务 1e6c3e2e，2026-08-28 定稿版）}：本查询共检出 31 项跨源冲突（26 项自动归因、0 项需人工），按（实体、字段）聚合做方差分析。以下为其中两组真实示例：

\begin{longtable}{@{}p{1.7cm}p{3.4cm}p{3.6cm}p{3.2cm}p{2.4cm}@{}}
\toprule
冲突字段 & 冲突对（来源 A vs 来源 B） & 取值与统计量 & 系统归因 & 实际处理 \\
\midrule
distance & MWSC 星团目录 vs 2018Ap\&SS 论文（20 个来源参与方差分析） & 130.0 pc vs 136.8 pc；Cohen's d = 2.29，95\% CI [$-$3.8, 8.4]，单位一致、方法与条件不同 & statistical\_outlier——目录年代久远的历史值 vs 现代测距值，方法/时代差异导致的统计离群 & 不删除、不合并：MWSC 130 pc 以独立记录保留并打离群标注，供研究者知悉其为旧目录值 \\
cluster\_age & MWSC 星团目录 vs 2018ApJ 论文（对数 vs 线性写法） & logt 8.15 vs 112.0 Myr（vs 2018A\&A 125.0 Myr 同型）；Cohen's d $>$ 10，单位口径不同 & 单位口径不一致（对数写法 log(yr) 与线性 Myr）形成的极端伪离群 & 对数写法按“无法可靠互算”保留原始写法并标注（见 3.4.4 记录状态表“待核对”行），不换算进 Myr 主流组 \\
\bottomrule
\end{longtable}

\noindent 31 项冲突的处理去向全部为“保留并标注”——系统只归因、不擅裁，取值分群与涉及来源写进每条记录的冲突标注，供研究者逐条回溯裁决（完整冲突清单随附录交付）。

\bigskip
\noindent\textbf{结构化输出中的来源关系：}

\begin{itemize}[leftmargin=1.2em,itemsep=2pt]
\item \textbf{每条记录实际保留了哪些来源信息}：论文记录保留文献编号、页码、原文定位框与原文摘录；星表记录保留数据表、键列、键值与行号；补充数据表记录保留数据表编号与行号；所有记录保留提取方式、提取置信度与原始写法。
\item \textbf{多个来源支持同一字段时如何表示}：同一物理量的多来源取值以独立记录并存，每条携带各自来源；取值分群时在冲突标注中列出涉及来源与值域区间，研究者可逐条回溯。
\item \textbf{模型补充、规则换算和原始数据如何区分}：原始数据（从原文逐字提取或星表查询的记录）与规则换算（单位统一等，留修改记录）由修改记录区分；模型生成的清洗与洞察输出一律标注为模型生成，不与原始数据混同。
\item \textbf{团队没有合并的冲突数据及原因}：不同文献对同一物理量的取值差异（如早期与近代测距结果）全部保留并标注方法差异——系统不擅自动用领域知识裁决，取舍交由研究者，以保证数据的可追溯性与科研判断的独立性。
\end{itemize}

\subsubsection{输出结果展示}

质量管线走完后的\textbf{记录终态一览}：下面以 M45 案例展示“原文 → 提取 → 统一/裁决 → 状态”的代表性输出；记录状态（正确/已统一/待核对）即 P10、P11 各环节处理后的真实结果，其中“待核对”项与第 4 章 V1 版本表（同批记录）呼应。

\noindent\textbf{M45 代表性输出一览：}

\begin{longtable}{@{}p{3.4cm}p{2.3cm}p{2.6cm}p{3.0cm}p{2.4cm}@{}}
\toprule
原始内容或位置 & 提取结果 & 单位或上下文 & 来源关系 & 当前状态（正确/待核对/冲突） \\
\midrule
1999ApJ 第 13 页正文：“Hipparcos 视差误差相关性”上下文 & 130.7 $\pm$ 11.1 & pc & 1999ApJ...523.328N，正文提取 & 正确 \\
2005A\&A 第 1 页表格（距离与金属丰度测定表） & $-$0.034 $\pm$ 0.024 & dex & 2005A\&A...429.645S，表格提取 & 正确 \\
2018MNRAS 表格单元格 “7.48 $\pm$ 0.03” & 7.48 & mas & 2018MNRAS.473.4612K，表格提取 & 已统一（清理不确定度，修改记录可回溯） \\
MWSC 目录列 logt & 8.15 & log(yr)（原始写法保留） & MWSC 星团目录，数据库查询 & 待核对（与论文 Myr 写法不同，保留原文标注） \\
2006A\&A 第 2 页 “星团参数汇总表” & 75-100 & Myr（区间写法保留原样） & 2006A\&A...446.611F，表格提取 & 待核对（区间写法按设计保留） \\
\bottomrule
\end{longtable}

\noindent\textbf{M45 收敛结果}：五项物理量经统一与冲突裁决后形成可对比的统计收敛——距离 37 条记录集中在 134.4--136.2 pc、距离模数 25 条集中在 5.58--5.65 mag、视差 7 条集中在 7.42--7.49 mas、年龄 26 条集中在 100--130 Myr、金属丰度 8 条集中在 $-$0.03~+0.03 dex；早期与现代测距的分歧作为真实冲突保留并标注（31 项，见 3.4.3）。全量输出与逐条状态随附录交付。

\bigskip
\noindent\textbf{结构化数据包（五件套 + 附加证据）}。每次查询产出的数据包构成：结构化数据文件、长表与宽表（CSV）、质量评分与质量报告、来源信息、修改记录与洞察报告，以及论文图表证据与原文定位页面；前端与导出同源，供研究者下载与核验。文件由导出链统一生成并校验：

\begin{longtable}{@{}p{2.4cm}p{4.2cm}p{7.8cm}@{}}
\toprule
数据包构成 & 实际文件 & 内容与用途 \\
\midrule
① 结构化记录 & \code{grounded\_data\_*.json} & 全部记录的结构化本体：数值/单位/字段（受性质清单约束）、来源、提取方式与置信度、原文摘录与测量方法标注——供程序直接读取 \\
② 长表 / ③ 宽表 & \code{data\_long\_*.csv}、\code{data\_wide\_*.csv} & 记录级长表与实体级宽表（含页码/定位框/置信度列），可直接导入分析脚本与绘图 \\
④ 质量报告 & \code{quality\_summary\_*.json} & 综合质量分（含等级与置信区间）、逐源得分、置信度分解、问题汇总与风险指标、使用建议 \\
⑤ 来源与溯源 & \code{metadata\_*.json} + \code{traceability\_*.json} & 字段字典、来源汇总、处理记录（含冲突 31/26 归因统计）；修改血缘（输入→输出记录映射、每次修改的前后值与原因）——“谁能回溯到哪”的完整凭证 \\
附加 & \code{insights\_*.json}、\code{manifest\_*.json} & 洞察报告（字段关联/使用建议/综合叙述，标注模型生成）；导出清单与校验结果（行数、校验是否通过） \\
\bottomrule
\end{longtable}

\noindent 另含论文图表证据（裁剪图 + 图注 + 定位）与原文定位页面，在前端图证画廊与溯源视图中呈现（见 3.5）。

\bigskip
\noindent\textbf{洞察报告}：在结构化数据交付后，系统生成字段关联分析、跨字段关系、使用建议与综合叙述报告，全部明确标注为模型生成（qwen3.7-flash），与原始数据区分，供研究者核验与使用。

\subsection{前端与交互（产品特色）}

\begin{figure}[H]\centering
\includegraphics[width=0.9\textwidth]{\detokenize{frontend.png}}
\caption{前端双分支功能：我的查询/演示样例双栈、增值面板与交付层}\label{fig:frontend}
\end{figure}

前端将“一次查询的一生”呈现为可交互的时间线：单页三栏（任务列表 / 对话与阶段卡 / 右侧详情面板），对话与工作流两种视图共享同一条 SSE 实时状态流，研究者全程可见、可点、可回放。主要交互特点：

\begin{enumerate}[leftmargin=1.6em,itemsep=2pt]
\item \textbf{全过程时间线}：消息与理解/检索/提取/质量检查/清洗/交付/洞察阶段卡按真实事件序交错推进，卡头带状态点、耗时与摘要，点击展开内部步骤；清洗卡内“规范化/冲突消解/人工审核”以第 N 轮嵌套块呈现，流程阶段一目了然。
\item \textbf{工作流视图与三级下钻}：同一查询可切换为编号圆点纵向时间线（各阶段耗时条归一化对比），逐级下钻到阶段明细（步骤/流程轮次/Agent/日志），再到 Agent 级修改轨迹——后端流程对研究者完全透明。
\item \textbf{澄清与人工复核（HITL）}：需要确认处澄清卡挂起逐轮提问；需要裁决处人工审核面板批量展示冲突（来源 A/B 对照、分歧原因、逐项 1--5 裁决 + 可选理由），强制逐项选择后一次提交，理由随修改记录进入最终导出溯源。
\item \textbf{记录表与血缘下钻}：结果以记录表原位呈现（实体/性质/值/单位/来源/提取方式），点任意一行展开单条血缘：星表坐标或论文页码与 bbox 定位缩略图、原文摘录、处理轨迹（原值→终值+修改步骤）、冲突标注与模型洞察。
\item \textbf{图证画廊与全屏核验}：论文图证据以缩略图双列呈现，点开全屏浏览——bbox 以三种样式区分数值定位/整表定位/兜底段落定位，支持光标锚定缩放与拖拽，供研究者逐张核对证据。
\item \textbf{质量与运行详情面板}：综合质量分大字 + 六维雷达、置信区间、逐源评分、路由分布与决策矩阵、多源方差与冲突标注；运行页列出各环节 LLM 与工具调用计数与审计时间线，模型参与程度可核验。
\item \textbf{双模式与全保真回放}：“我的查询 ⇄ 演示样例”一键切换；17 组内置样例均为真实运行任务（只读本体、副本可删），历史任务按压缩节奏全保真重演、不消耗调用额度，底部日志抽屉逐行审计——评审与复现无需重跑系统。
\item \textbf{导出与数据包}：性质勾选 + 来源分组 + 记录行级勾选预览，导出的长表 CSV（含页码/定位框/置信度列）或完整 JSON 与看板展示同源；数据包五件套本地保存，供下游分析直接使用。
\end{enumerate}

\noindent\emph{（界面截图以主界面时间线全景与演示样例画廊两张为主，随提交材料附上；细节交互见上条文字特点）}
